\section{Background}

Recursive Spawning addresses an emerging class of problems at the intersection of multi-agent systems, distributed accountability, and hierarchical task decomposition. To situate the contribution precisely, this section reviews the prior art across four interdependent domains: agent lifecycle management in multi-agent systems (§\ref{sec:agent-lifecycle}), accountability and provenance in distributed systems (§\ref{sec:provenance}), fixed versus mutable delegation structures (§\ref{sec:delegation-chains}), and hierarchical task decomposition in AI (§\ref{sec:task-decomp}). We conclude by positioning the BX3 Framework as a concrete architectural substrate that resolves tensions across these threads through immutable parent chain semantics (§\ref{sec:bx3-substrate}).

\subsection{Agent Lifecycle Management in Multi-Agent Systems}
\label{sec:agent-lifecycle}

The operational lifecycle of autonomous agents in distributed AI systems presents compounding governance challenges that scale non-linearly with agent count and interaction depth. Modern multi-agent systems (MAS) require lifecycle management frameworks that address provisioning, role-binding, state synchronization, and graceful retirement of agent instances \cite{cenfic_production_2025, lumenova_govern_2025}. Unlike monolithic single-agent deployments, MAS lifecycles involve nested spawn chains, emergent inter-agent dependencies, and coordination states that are difficult to audit post-hoc.

Agent lifecycle management in production environments encompasses five primary phases: \textbf{Provisioning} — the instantiation and role-binding of a new agent; \textbf{Activation} — the transition from idle to task-active state; \textbf{Delegation} — the assignment of authority and task scope; \textbf{Monitoring} — continuous state tracking and health verification; and \textbf{Retirement} — graceful termination with state archival and audit closure. Current industry frameworks handle these phases inconsistently: hierarchical orchestrators (e.g., manager-worker patterns) manage lifecycle top-down but create single points of failure; peer-to-peer agent meshes distribute authority but introduce provenance gaps where actions cannot be traced to a single accountability anchor \cite{voltagent_papers_2026}.

Critically, most existing lifecycle frameworks treat the delegation chain as a transient, mutable construct — permissions are granted, modified, or revoked during runtime. This creates what the BX3 Framework identifies as a \textbf{Logic Collision} risk: reasoning and execution occupy the same functional plane, enabling actions to escape deterministic verification. The SentinelAgent framework formalizes this as the \textbf{delegation accountability gap}: in a chain where Agent A delegates to Agent B, which invokes Tool C on behalf of User X, no existing framework can definitively answer \emph{whose authorization chain led to this action, and where did it violate policy?} \cite{sentinel_agent_2026}. The Delegation Chain Calculus (DCC) introduced in that work models each delegation step as a token carrying authority scope, intent vector, and policy constraints — formalizing seven properties including authority monotonic narrowing, intent preservation, and forensic reconstructibility. However, even the DCC acknowledges that intent verification (Property P2) is \emph{probabilistic} rather than deterministic, due to the fundamental incomputability of semantic intent classification over natural language delegation descriptions \cite[Proposition 1]{sentinel_agent_2026}.

The BX3 Framework's Loop Isolation principle architecturally separates Purpose (accountability anchor), Bounds Engine (propositional logic), and Fact Layer (physical execution) into discrete planes, making the delegation chain a structural property rather than a runtime convention. This design choice directly addresses the mutable-chain failure mode by making parent-child relationships immutable by construction.

\subsection{Accountability and Provenance in Distributed Systems}
\label{sec:provenance}

Distributed AI systems face a fundamental accountability deficit: actions performed by autonomous agents generate outcomes that are computationally complex to attribute, operationally difficult to audit, and legally ambiguous to assign liability for. This deficit is amplified when agents spawn sub-agents that perform further actions — each layer of delegation dilutes the traceability of the original authorization. The operational complexity of multi-agent systems can demand up to 26 times the monitoring resources required for single-agent deployments \cite{lumenova_govern_2025}, creating an auditing burden that scales superlinearly with chain depth.

The W3C PROV standard provides a provenance model for tracking the derivation history of data products, but its application to AI agent actions — including tool calls, LLM invocations, prompt modifications, and workflow state transitions — requires substantial extension. PROV-AGENT addresses this gap by extending the PROV model with AI-specific artifacts (tools, prompts, responses, model invocations) unified in a single provenance graph \cite{prov_agent_2025}. This enables tracing erroneous outputs back to upstream prompts, inputs, and prior agent decisions — a capability that existing distributed AI monitoring tools largely lack.

Beyond provenance tracking, accountability in distributed systems requires cryptographically sealed audit logs that survive agent retirement and organizational restructuring. The LLM Delegate Protocol (LDP) extension work identified a \textbf{provenance paradox}: when delegates can inflate self-reported quality scores, quality-based routing systematically selects the worst performers, performing worse than random selection (simulated: 0.55 vs. 0.68 mean output quality; real-model validation confirmed the same inversion with degraded Claude Haiku capturing 100\% of routing despite lowest true quality) \cite{provance_paradox_2026}. This finding demonstrates that provenance metadata — when self-reported rather than independently attested — can actively degrade system outcomes rather than improve them. The LDP extensions introduce \textbf{delegation contracts} that bound authority through explicit objectives, budgets, and failure policies; a \textbf{claimed-vs-attested identity model} distinguishing self-reported from verified quality; and \textbf{typed failure semantics} enabling automated recovery.

The BX3 Framework's forensic ledger approach (9-Plane DAP architecture) addresses this by maintaining Plane-7 (Authorization Ledger) and Plane-8 (Financial Ledger) as immutable, independently loggable records for every event across Purpose, Bounds Engine, and Fact Layer planes \cite{bx3_universal_spec_2026}. Each plane maintains its own tamper-evident log chain such that compromise of any single plane does not propagate to others — a property absent from most centralized monitoring systems where a single compromised audit store invalidates the entire audit history. The Training Wheels Protocol introduces a complementary behavioral accountability layer by governing not only what agents do, but under what conditions they are permitted to act autonomously versus requiring human review \cite{training_wheels_protocol_2026}. Under Training Wheels, every outbound action is queued for human review before execution until the agent earns autonomy through sustained approved behavior — creating a retrospective accountability record qualitatively different from pre-action provenance graphs alone.

\subsection{Fixed vs Mutable Delegation Chains}
\label{sec:delegation-chains}

Delegation chains in multi-agent systems can be classified along a spectrum of mutability. \textbf{Mutable delegation chains} — where permissions, scopes, and parent-child relationships are modified dynamically during runtime — offer flexibility but introduce dangerous ambiguity: when an agent acts, it may be unclear whether its authorization originated from its creator, from a mid-chain modification, or from an inheritance pattern that no longer reflects the current organizational state \cite{sentinel_agent_2026}.

\textbf{Fixed delegation chains} — where the parent-child relationship and its associated permissions are immutable once established — provide stronger provenance guarantees at the cost of reduced runtime adaptability. Static delegation is common in hierarchical task decomposition systems where a planner agent decomposes a task into sub-tasks assigned to fixed worker roles \cite{claude_code_design_2026, openai_mas_2025}. The worker agents inherit their task scope from the parent planner's decomposition and cannot expand their own authority without a new spawn event. Google's Agent-to-Agent (A2A) Protocol and Anthropic's Model Context Protocol (MCP) provide skill-based routing and tool integration at the communication layer \cite{a2a_protocol_2025, mcp_protocol_2024}, but neither protocol-level mechanism governs \emph{what happens when a delegated task must itself be further delegated} — a recursive delegation scenario that the BX3 Framework's Recursive Spawning directly addresses.

The distinction is architecturally significant. In mutable systems, a mid-chain agent can modify its own permissions or spawn unauthorized children without a top-down authorization event, creating a provenance gap. In the SentinelAgent threat model, this corresponds to Adversary Type A4 (Privilege Escalator) and A5 (Chain Poisoner) \cite[Table I]{sentinel_agent_2026}. The DCC addresses this through Property P1 (Authority Monotonic Narrowing) — enforcing that $\sigma_{i+1} \subseteq \sigma_i$ at every delegation step — and Property P4 (Forensic Reconstructibility) — maintaining a hash-linked chain where each token $\tau_i$ contains $h_p = H(\tau_{i-1})$, enabling $O(n)$ reconstruction of the full chain for audit. However, the DCC's chain is mutable in the sense that the DAS can issue revised tokens; the immutability is in the historical record, not in the protocol's resistance to mid-chain reauthorization.

The BX3 Framework's position on this spectrum is explicit: delegation chains are \textbf{structurally fixed} at spawn time but \textbf{hierarchically updatable} only via a controlled Recursive Spawning event. A Parent Node births a Child Loop by provisioning a Worksheet — a containerized, self-contained logic set delivered Over-the-Air (OTA) to the target node — that includes a hard-coded pointer to the Parent's Purpose layer, preventing autonomous drift while enabling local survivability during cloud disconnection \cite{bx3_universal_spec_2026}. The delegation is fixed in the sense that the child cannot alter its parent pointer after spawn; it is updatable only through a new Recursive Spawning event initiated by the parent, which creates a new Worksheet with updated parameters. This design prevents the mid-chain mutation failure mode: any modification to the delegation structure requires a top-down Recursive Spawning event, ensuring the Human Root anchor at Level 0 maintains a verifiable, non-collapsing chain of delegation to every active node.

\subsection{Hierarchical Task Decomposition in AI}
\label{sec:task-decomp}

Hierarchical task decomposition (HTD) is a foundational pattern in autonomous AI systems, enabling complex goals to be recursively broken into manageable sub-tasks assigned to specialized agents. Manager-worker architectures decompose high-level objectives into a tree of sub-tasks, where a planner agent generates task hierarchies and delegates execution to spawned sub-agents with scoped tool permissions and context isolation \cite{lumenova_govern_2025, sap_agent_arch_2025}.

HTD systems face three structural challenges as they scale: \textbf{decomposition fidelity} — whether sub-task granularity accurately reflects the true causal structure of the parent task; \textbf{delegation auditability} — whether each sub-agent's authorization and execution is traceable back to the original decomposer; and \textbf{error propagation} — whether failures in sub-agents cascade upward with sufficient context for the parent to diagnose and recover without human intervention. Research in AI agent design identifies hierarchical orchestration as a primary mechanism for managing complexity in multi-agent workflows \cite{linkedin_ai_agents_2025, futran_modelops_2025}. Claude Code's architecture exemplifies this pattern: the parent spawns sub-agents (Explore, Plan, general-purpose, and custom types) that operate with restricted tool sets and return summary-only results, creating a delegation hierarchy that is structurally clean but operationally opaque without additional instrumentation \cite{claude_code_design_2026}.

A key limitation in most HTD implementations is that spawned agents typically receive a task description without inheriting the parent's governance constraints. This creates a structural gap: the parent agent's intent, scope, and policy boundaries are communicated as runtime context rather than encoded as architectural properties of the child. Consequently, a deeply spawned agent can lose coherent alignment with the original intent, or accountability becomes diffuse across multiple levels of delegation. The BX3 Framework's Recursive Spawning mechanism addresses this by embedding the decomposition context — the parent's Purpose layer and operational bounds — directly into the Child Loop's Worksheet at spawn time, making the parent pointer and bounds properties of the spawned agent's structural definition rather than its runtime context.

\subsection{The BX3 Framework as Substrate for Immutable Parent Chains}
\label{sec:bx3-substrate}

The BX3 Framework provides a structural substrate in which immutable parent chains are not a policy choice but an architectural consequence. Three mutually reinforcing mechanisms enforce this property: \textbf{Loop Isolation}, \textbf{Recursive Spawning}, and \textbf{Bailout Protocol}.

\textbf{Loop Isolation} ensures that every node — from the Human Root at Level 0 to the most distant field sensor — is a self-contained BX3 Loop consisting of Purpose, Bounds Engine, and Fact Layer, strictly separated into discrete planes. A reasoning engine (Bounds Engine) never shares a plane with the physical execution layer (Fact Layer), making Logic Collisions architecturally impossible regardless of the agent's autonomy level \cite{bx3_universal_spec_2026}. This is a stronger guarantee than DCC's Property P1 (authority narrowing), which constrains delegation tokens but does not architecturally separate reasoning from execution within an agent.

\textbf{Recursive Spawning} operationalizes immutable parent chains by defining the spawn event as a one-way provisioning operation: a Parent Node births a Child Loop by deploying a Worksheet containing complete Bounds Engine and Fact Layer logic for the target environment, along with a hard-coded, immutable pointer to the Parent's Purpose layer. This pointer is not a reference that can be overwritten at runtime — it is a structural property encoded in the Worksheet's foundation. The child loop inherits its operational scope entirely from the parent and cannot re-parent itself without a new Recursive Spawning event initiated by the parent or an ancestor. This contrasts with mutable delegation systems where mid-chain agents can autonomously expand their own scope or spawn unauthorized children \cite{sentinel_agent_2026}.

The \textbf{Bailout Protocol} ensures that immutable parent chains do not become liability traps. When a sub-node encounters a condition outside its provisioned operational range, it propagates an exception asynchronously up the recursive tree, bypassing all Machine actors, until anchored by a Human Accountability Anchor. The trigger lifecycle (TRIGGER\_FIRED $\rightarrow$ DELIVERING\_TO\_PARENT $\rightarrow$ PARENT\_EVALUATING $\rightarrow$ ESCALATING\_TO\_PARENT $\rightarrow$ HUMAN\_ROOT\_REVIEWING) ensures that every unresolved exception eventually reaches a human with appropriate authority — maintaining chain integrity without creating a deadlock condition \cite{bx3_universal_spec_2026}.

The 9-Plane DAP architecture underpinning the BX3 Framework provides the logging substrate for chain integrity. Each plane (from Plane-1 Physical Substrate through Plane-9 IP Core) maintains independently auditable records, and the Spatial Firewall enforces plane isolation such that data on Plane N cannot be used to infer or reconstruct data on Plane N+1 or higher \cite{bx3_universal_arch_2026}. This ensures that even if a compromised node generates false provenance data, the evidence is confined to its own plane and cannot contaminate the chain of custody. The combination of structurally immutable parent pointers, plane-isolated forensic logs, and asynchronous Bailout escalation produces an accountability architecture that is qualitatively distinct from prior art: every descendant agent is anchored to exactly one Human Root by architectural definition, not by convention or policy.

% Bibliography
\bibliographystyle{plain}
\bibliography{bx3framework}
